% =====================================================================
%  CARNET D'ENTRAÎNEMENT — Fiche 2 : Calcul vectoriel
%  Thème 5 — Angle entre deux vecteurs et identités remarquables — TUTORIEL
% =====================================================================
\documentclass[11pt,a4paper]{article}
\usepackage[utf8]{inputenc}
\usepackage[T1]{fontenc}
\usepackage{lmodern}
\usepackage[french]{babel}
\usepackage{amsmath,amssymb,mathtools}
\usepackage{xcolor}
\usepackage{colortbl}
\usepackage{tikz}
\usepackage[most]{tcolorbox}
\usepackage{enumitem}
\usepackage{booktabs}
\usepackage{fancyhdr}
\usepackage[a4paper,margin=2.1cm,headheight=26pt,headsep=10pt]{geometry}
\usetikzlibrary{arrows.meta,positioning,calc,angles,quotes}

\definecolor{primary}{RGB}{222,70,80}
\definecolor{primarydark}{RGB}{150,30,42}
\definecolor{methcol}{RGB}{90,90,100}
\definecolor{excol}{RGB}{0,135,90}
\definecolor{ideecol}{RGB}{120,94,200}
\definecolor{formulacol}{RGB}{198,93,9}
\definecolor{attncol}{RGB}{200,40,55}
\definecolor{vecu}{RGB}{20,20,20}
\definecolor{vecv}{RGB}{0,110,200}
\definecolor{axiscol}{RGB}{90,90,90}

\tcbset{boxbase/.style={enhanced,breakable,boxrule=0.9pt,arc=2.5pt,
  left=3mm,right=3mm,top=2.2mm,bottom=2.2mm,fonttitle=\bfseries,coltitle=white}}
\newtcolorbox{methode}[1][]{boxbase,colback=methcol!7,colframe=methcol,sharp corners,
  title={\textsc{Comment faire ?}\ifx\relax#1\relax\relax\else\textmd{ --- \itshape #1}\fi}}
\newtcolorbox{exemple}[1][]{boxbase,colback=excol!5,colframe=excol,
  title={\textsc{Exemple}\ifx\relax#1\relax\relax\else\textmd{ : \itshape #1}\fi}}
\newtcolorbox{idee}[1][]{boxbase,colback=ideecol!6,colframe=ideecol,
  title={\textsc{Intuition}\ifx\relax#1\relax\relax\else\textmd{ : \itshape #1}\fi}}
\newtcolorbox{rappel}[1][]{boxbase,colback=formulacol!7,colframe=formulacol,
  title={\textsc{À retenir}\ifx\relax#1\relax\relax\else\textmd{ --- \itshape #1}\fi}}
\newtcolorbox{piege}[1][]{boxbase,colback=attncol!6,colframe=attncol,
  title={\textsc{Attention}\ifx\relax#1\relax\relax\else\textmd{ : \itshape #1}\fi}}

\pagestyle{fancy}\fancyhf{}
\fancyhead[L]{\small\textcolor{primarydark}{\textbf{Fiche 2} — Calcul vectoriel}}
\fancyhead[R]{\small\textcolor{primarydark}{\textsc{Thème 5 — Tutoriel}}}
\fancyfoot[C]{\small\thepage}
\renewcommand{\headrulewidth}{0.8pt}

\newcommand{\vc}[1]{\overrightarrow{#1}}
\providecommand{\norm}[1]{\left\lVert #1\right\rVert}

\begin{document}
\begin{center}
{\LARGE\bfseries\color{primarydark}Thème 5 — Angle entre deux vecteurs et identités remarquables}\\[2pt]
{\color{primary}\rule{0.82\linewidth}{1.2pt}}\\[4pt]
{\small\itshape Du produit scalaire à l'angle : mesurer, et manipuler les normes de sommes.}
\end{center}
\vspace{1mm}

\begin{rappel}[cosinus de l'angle entre deux vecteurs]
\[
  \boxed{\;\cos(\vec u,\vec v)=\dfrac{\vec u\cdot\vec v}{\norm{\vec u}\,\norm{\vec v}}
  =\dfrac{x x'+y y'}{\sqrt{x^2+y^2}\,\sqrt{x'^2+y'^2}}\;}
\]
On calcule le produit scalaire et les deux normes, puis $\theta=\arccos(\dots)$. Le cosinus est un
nombre de $[-1;1]$ : positif $\to$ angle aigu, nul $\to$ droit, négatif $\to$ obtus.
\end{rappel}

\begin{idee}[reconnaître les angles « remarquables »]
Certaines valeurs du cosinus donnent des angles à connaître par cœur — inutile de sortir la
calculatrice :
\end{idee}

\begin{center}
\renewcommand{\arraystretch}{1.3}
\begin{tabular}{ccccc}
\toprule
\rowcolor{primary!12}$\cos\theta$ & $0$ & $\dfrac12$ & $\dfrac{\sqrt2}{2}$ & $\dfrac{\sqrt3}{2}$\\
\midrule
$\theta$ & $90^\circ$ & $60^\circ$ & $45^\circ$ & $30^\circ$\\
\bottomrule
\end{tabular}
\end{center}

\begin{methode}[trouver l'angle entre $\vec u$ et $\vec v$]
\begin{enumerate}[leftmargin=1.7em,topsep=2pt,itemsep=2pt]
  \item Calculer le produit scalaire $\vec u\cdot\vec v=xx'+yy'$.
  \item Calculer les deux normes $\norm{\vec u}$ et $\norm{\vec v}$.
  \item Former le quotient $\cos\theta=\dfrac{\vec u\cdot\vec v}{\norm{\vec u}\,\norm{\vec v}}$
        (le \emph{simplifier}).
  \item Reconnaître un angle remarquable, ou prendre $\theta=\arccos(\dots)$ à la calculatrice.
\end{enumerate}
\end{methode}

\begin{exemple}[un angle remarquable]
$\vec u\,(0;2)$ et $\vec v\,(2;2\sqrt3)$. On a $\vec u\cdot\vec v=0\times2+2\times2\sqrt3=4\sqrt3$,
$\norm{\vec u}=2$ et $\norm{\vec v}=\sqrt{4+12}=4$. Donc
\[
  \cos\theta=\frac{4\sqrt3}{2\times4}=\frac{\sqrt3}{2}\ \Longrightarrow\ \theta=30^\circ.
\]
\end{exemple}

\begin{rappel}[identités remarquables vectorielles]
\[
  \boxed{\;\norm{\vec u+\vec v}^2=\norm{\vec u}^2+\norm{\vec v}^2+2\,\vec u\cdot\vec v\;}
  \qquad
  \boxed{\;\norm{\vec u-\vec v}^2=\norm{\vec u}^2+\norm{\vec v}^2-2\,\vec u\cdot\vec v\;}
\]
\[
  \boxed{\;(\vec u+\vec v)\cdot(\vec u-\vec v)=\norm{\vec u}^2-\norm{\vec v}^2\;}
\]
Ce sont les analogues de $(a\pm b)^2=a^2+b^2\pm2ab$ et $(a+b)(a-b)=a^2-b^2$, le terme
$\pm2\,\vec u\cdot\vec v$ jouant le rôle du « double produit ».
\end{rappel}

\begin{exemple}[deux façons de calculer $\norm{\vec u+\vec v}$]
$\vec u\,(1;1)$, $\vec v\,(2;3)$. \emph{Méthode directe} : $\vec u+\vec v\,(3;4)$, donc
$\norm{\vec u+\vec v}=\sqrt{9+16}=5$.\\
\emph{Par l'identité} : $\norm{\vec u}^2=2$, $\norm{\vec v}^2=13$, $\vec u\cdot\vec v=2+3=5$, donc
$\norm{\vec u+\vec v}^2=2+13+2\times5=25$, soit $5$. Même résultat.
\end{exemple}

\begin{piege}[erreurs à éviter]
\begin{itemize}[leftmargin=1.5em,topsep=2pt,itemsep=1pt]
  \item $\norm{\vec u+\vec v}\neq\norm{\vec u}+\norm{\vec v}$ en général ! Il y a le double produit
        $2\,\vec u\cdot\vec v$.
  \item Bien mettre les normes (jamais les carrés) au dénominateur du cosinus.
  \item Le cosinus doit tomber dans $[-1;1]$ : sinon, il y a une erreur de calcul.
\end{itemize}
\end{piege}

\end{document}
